% !TEX program = pdflatex
\documentclass[12pt,a4paper]{letter}

\usepackage[T1]{fontenc}
\usepackage[margin=1in]{geometry}
\usepackage{hyperref}
\usepackage{setspace}
\usepackage{parskip}

\hypersetup{
  colorlinks=true,
  urlcolor=blue
}

\onehalfspacing

\signature{Professor Ayo Ogunjuyigbe\\
Professor of Electrical \& Electronic Engineering\\
Former Director, Management Information Systems (MIS)\\
University of Ibadan, Nigeria}

\address{Department of Electrical \& Electronic Engineering\\
University of Ibadan\\
Ibadan, Oyo State\\
Nigeria}

\date{27 February 2026}

\begin{document}

\begin{letter}{UK Visas \& Immigration\\
Global Talent Visa --- Tech Nation (Digital Technology)\\
United Kingdom}

\opening{To Whom It May Concern,}

I write to provide my personal endorsement of Odeyemi Olusegun Israel in support of his application for the UK Global Talent Visa under the Digital Technology route. While my colleague Professor {[CS Professor Name]} is better placed to comment on the theoretical merits of Olusegun's research papers, my endorsement focuses on what I consider equally important: his \textbf{engineering capability, systems thinking, and the real-world impact of the products he has built}.

By way of background, I am a Professor of Electrical and Electronic Engineering at the University of Ibadan and have previously served as Director of Management Information Systems (MIS) at the same institution. In the MIS role, I was responsible for overseeing the university's entire digital infrastructure---enterprise systems, data management, and operational technology. My research spans power systems, computational intelligence, and applied machine learning in engineering contexts. This combination of institutional systems management and technical research gives me a particular perspective on what it takes to build complex, production-grade systems that actually work.

\textbf{AMTTP: a product, not a prototype.}

Olusegun has designed and built a system called the Anti-Money Laundering Transaction Trust Protocol (AMTTP). I want to be clear about what this is, because in my experience, most solo developers who claim to have built ``platforms'' have really built proof-of-concept demos. AMTTP is not that.

The system comprises four architectural layers that must function together in real time:

\begin{enumerate}
  \item A \textbf{client SDK layer} (TypeScript and Python, both open-sourced on GitHub) that financial institutions integrate into their transaction pipelines.
  \item A \textbf{compliance orchestration engine} that combines ML risk scoring, graph-based relationship analysis, sanctions list screening, and configurable policy rules into a single deterministic pass/fail decision.
  \item An \textbf{offline ML training pipeline} using a Composite Teacher architecture: Autoencoder-Enhanced XGBoost (weighted at 0.4), seven FATF-aligned AML pattern detectors (0.3), and graph structural properties (0.3) to generate pseudo-labels across 2.64 million transactions. A Student model distils this into a deployable classifier.
  \item A \textbf{blockchain smart contract layer} with 18 contracts deployed on Ethereum Sepolia, handling on-chain risk attestation, threshold-based transaction gating, and compliance audit trails.
\end{enumerate}

The whole thing runs across 17 containerised microservices coordinated through a central API gateway. I know, from my time managing the University of Ibadan's IT systems, how difficult it is to keep even five or six services reliably communicating. Seventeen is a fundamentally harder problem, and Olusegun has solved it alone.

\textbf{Why the architecture matters.}

The core insight behind AMTTP addresses a problem that the entire blockchain compliance industry has, to date, failed to solve. In traditional banking, a suspicious transaction can be frozen mid-flight. On a blockchain, once a transaction reaches finality, it is irreversible. Every major compliance tool currently on the market---Chainalysis, Elliptic, TRM Labs---operates \textit{after} settlement. They generate alerts; they do not prevent illicit transactions from completing. AMTTP intervenes \textit{before} settlement, at the protocol level.

This is not simply a matter of calling an API earlier in the pipeline. It requires solving a hard real-time systems problem: the ML inference, graph analysis, sanctions screening, and policy evaluation must all complete within the block confirmation window (approximately 12 seconds on Ethereum). Olusegun has engineered the system to meet these timing constraints while maintaining deterministic compliance guarantees. That is genuine systems engineering, not application development.

\textbf{Zero-knowledge cryptography for regulatory compliance.}

One component of AMTTP that I think deserves particular attention is zkNAF---Olusegun's zero-knowledge proof framework for privacy-preserving compliance verification. It allows institutions to prove facts about a customer (KYC validity, risk score within acceptable range, absence from sanctions lists) without revealing the underlying personal data.

From a regulatory perspective, this is significant. Financial institutions currently face a genuine conflict between GDPR's data minimisation requirements and AML regulations' transparency demands. Olusegun's solution does not merely acknowledge this tension---it resolves it architecturally. He has also implemented multi-oracle threshold signatures and replay protection, which are the kind of infrastructure-level security measures that separate production systems from academic exercises.

\textbf{Interdisciplinary range and engineering maturity.}

What distinguishes Olusegun, in my observation, is not just depth in any one area but the ability to work across domains at a high level simultaneously. AMTTP requires competence in:

\begin{itemize}
  \item \textbf{Machine learning engineering}: training pipelines, model deployment, feature engineering at scale.
  \item \textbf{Blockchain development}: Solidity smart contracts, gas optimisation, on-chain/off-chain coordination.
  \item \textbf{Distributed systems}: microservice orchestration, containerisation, API gateway design.
  \item \textbf{Cryptographic engineering}: zero-knowledge proofs, threshold signatures, secure key management.
  \item \textbf{Regulatory knowledge}: FATF Travel Rule, FCA requirements, GDPR data handling obligations.
\end{itemize}

In the university context, a project of this scope would typically require a principal investigator, two or three postdoctoral researchers, and a team of PhD students, with institutional support for infrastructure and compute. Olusegun has done it as a single engineer. I can confirm this because I have reviewed the GitHub repository and the deployment configuration---there is no evidence of collaborators, outsourced components, or wrapper libraries pretending to be original work. It is all his.

The decision to open-source the client SDKs while keeping the orchestration core proprietary also reflects good product judgment. It lowers the integration barrier for potential adopters while protecting the core intellectual property. This is the kind of strategic decision that separates engineers who can build things from engineers who can build \textit{products}.

\textbf{Relevance to the UK digital economy.}

The UK is actively positioning itself as a global leader in responsible fintech regulation. The FCA's forthcoming framework for crypto-asset regulation, expected by late 2027, will create substantial demand for compliance infrastructure that can operate at the protocol level. AMTTP is, as far as I am aware, the only system in existence that provides pre-settlement AML enforcement on blockchain transactions. The individual who built it would be a direct asset to the UK's regulatory technology sector.

Beyond AMTTP, Olusegun has also produced a body of research work (Blind Spot Decomposition Theory and the Universal Deviation Law) that has theoretical merit---but I defer to my colleague's letter for detailed comment on that aspect. What I will say is that the fact Olusegun can move fluently between building production systems and producing original mathematical theory is itself unusual and valuable.

\textbf{My assessment.} In my years at the University of Ibadan, both as a professor and as the director responsible for institutional IT infrastructure, I have mentored hundreds of engineering students and professionals. I can say without hesitation that Olusegun's engineering capability and systems-level thinking place him in the top tier of talent I have encountered. He does not simply use technologies; he designs and builds complex systems from first principles that solve problems no one else has solved yet.

I recommend him for the Global Talent Visa without reservation.

\closing{Yours faithfully,}

\end{letter}

\end{document}
